\textbf{17.} Докажите эквивалентность следующих двух определений логарифмической вычислимости функции $f\colon \{0, 1\}^* \to \{0, 1\}^*$:
\begin{itemize}
    \item[(a)] Длина $f(x)$ ограничена некоторым полиномом от длины $x$, а множества $A_s = \{(x, i)\colon i \le |f(x)|\}$ и $A_1 = \{(x, i)\colon f(x)_i = 1\}$ можно распознать на логарифмической памяти;
    \item[(b)] Функция $f$ вычисляется машиной с тремя лентами: входной, выходной и рабочей, при этом ответ записывается на выходной ленте, по которой можно двигаться только слева направо, а на рабочей ленте использовано логарифмическое число ячеек.
\end{itemize}

\begin{proof}

$(a)\Rightarrow(b)$ 

Мы знаем, что $|f(x)|\le p(|x|)$. Сделаем цикл $i=1\ldots p(|x|)$ (такой счётчик будет занимать $\log{p(|x|)} = O(\log{|x|})$ памяти), и в случае $(x,i)\in A_s$, проверим, что $(x,i)\in A_1$. В случае принадлежности вернем $1$, иначе $0$. Если же $(x,i)\not\in A_s$, то выведем символ конца вычислений и завершим работу. 

Так как распознание $A_s$ и $A_1$ работает на логарифмической памяти, то весь алгоритм можно сделать на логарифмической памяти.


$(b)\Rightarrow(a)$

Длина $f(x)$ ограничена $2\cdot$(число конфигураций), так как иначе мы бы просто зациклились. Следовательно, $|f(x)|\le p(|x|)$.

Для проверки $(x,i)\in A_s$ запустим вычисление $f$, но вместо записывания на выходную ленту будет увеличивать переменную длины, если она станет $\ge i$, то вернем $1$, в случае если мы выведем символ конца вычисления раньше, то вернем $0$.

Проверку $(x,i)\in A_1$ делаем аналогично, просто при записи $i$-го бита проверяем, что пишется $1$, в этом случае возвращаем $1$. Если же пишется $0$, или мы уже написали символ конца вычисления, то возвращаем $0$.

\end{proof}